\documentclass[12pt,a4paper]{article}
\usepackage[utf8]{inputenc}
\usepackage[T2A]{fontenc}
\usepackage[russian]{babel}
\usepackage{amsmath,amssymb,array}
\usepackage[margin=2.5cm]{geometry}
\usepackage[unicode,hidelinks]{hyperref}
\setlength{\emergencystretch}{3em}
\newcommand{\R}{\mathbb R}
\renewcommand{\C}{\mathbb C}
\title{Спектральная задача с отражением аргумента:\\
что известно и что получено}
\author{Учебный исследовательский проект}
\date{10 октября 2026 г.}
\begin{document}
\sloppy
\maketitle
\begin{abstract}
Обзор посвящён оператору второго порядка, в котором вторая производная в отражённой
точке входит с переменным коэффициентом. Мы объясняем, чем задача отличается от
изученных, какие теоремы из литературы к ней применимы, какие достаточные условия
базисности получены в проекте и что остаётся открытым. Доказательства изложены в
\texttt{exposition.tex}.
\end{abstract}

\section{Задача}
В пространстве $L_2(-1,1)$ рассматривается оператор
\[
 L_ay=-y''(x)-a(x)\,y''(-x),\qquad y(-1)=y(1)=0,
\]
где $y''(-x)$ --- вторая производная в точке $-x$. Коэффициент $a$ веществен, и при
некотором $\kappa<1$ почти всюду
\begin{equation}\label{eq:bound}
 |a(x)\,a(-x)|\le\kappa .
\end{equation}
В отличие от исходного черновика \texttt{main.tex}, произведение $a(x)a(-x)$ может
обращаться в нуль и быть отрицательным.

Вопрос: можно ли раскладывать функции по собственным и присоединённым (корневым)
функциям $L_a$ и насколько хорошо сходятся такие разложения. Различаются два ответа.
\emph{Базис Рисса} --- система, которая получается из ортонормированного базиса
ограниченным обратимым оператором; разложения по ней сходятся при любой перестановке
членов. \emph{Базис Рисса со скобками} слабее: корневые функции объединены в конечные
группы, и безусловно сходится ряд из сумм по группам. Второе свойство не влечёт первого,
и размер групп заранее неизвестен.

Трудность в том, что отражение стоит при старшей производной. Если $(Rf)(x)=f(-x)$, то
$L_a=(I+aR)(-d^2/dx^2)$, и при непостоянном $a$ оператор не самосопряжён и не является
возмущением оператора Дирихле младшими членами.

\section{Что известно}
\textbf{Операторы с отражением.} Для постоянного коэффициента собственные функции
выписываются явно и образуют ортонормированный базис~\cite{S2} (теорема~4.1, с.~4).
Бондаренко~\cite{S1} (\S2.1, с.~4) сводит уравнение с постоянной старшей частью к системе
для пары $(y(x),y(-x))$; этот приём мы используем как образец. Поляков~\cite{S5}
(теорема~1.1) методом подобных операторов доказывает базисность Бари для операторов
порядка выше второго. Работа~\cite{S6} относится ко второму порядку, но с другими
краевыми условиями. Во всех четырёх работах коэффициент при отражённой старшей производной
постоянен.

\textbf{Подчинённые возмущения.} Теорема Маркуса--Мацаева в формах~\cite{S4}
(теорема~6.1, с.~21) и~\cite{S3} (теорема~6.6, с.~140) утверждает следующее. Пусть
нормальный оператор имеет компактную резольвенту и спектр на конечном числе лучей,
возмущение подчинено ему с показателем $p<1$, а собственные значения расположены
достаточно редко. Тогда у возмущённого оператора есть базис Рисса из конечномерных
инвариантных подпространств, то есть базис со скобками. Снять скобки можно, если
константы Рисса внутри подпространств равномерны~\cite{S4} (предложение~3.2, с.~5);
достаточные условия требуют спектральных зазоров~\cite{S4} (теорема~6.2 и замечание~6.3,
с.~23--24). К этой же теме относятся асимптотика собственных значений возмущённых
операторов~\cite{S20} и локальная подчинённость форм~\cite{S11,S12}.

\textbf{Системы первого порядка.} Для систем типа Дирака строго регулярные краевые
условия дают базис Рисса, а регулярные --- базис Рисса со скобками~\cite{S8}
(теорема~1.1, с.~4; предложение~6.9, с.~34); см. также~\cite{S7,S9,S10}. Свести $L_a$ к
такой системе с областью, не зависящей от спектрального параметра, не удалось, поэтому
эти теоремы мы не применяем.

\textbf{От групп к отдельным функциям.} Для операторов Хилла и Дирака условия, при
которых скобки снимаются, изучены в~\cite{S13,S14,S15,S16,S17,S18,S21}; подход через
группы операторов дан в~\cite{S19}. Эти работы показывают, что нужны зазоры в спектре и
оценки спектральных проекторов.

Из работ~\cite{S7} и~\cite{S9}--\cite{S21} взяты только постановки задач; их теоремы мы не
проверяли и в доказательствах не используем. Номера теорем и страниц остальных работ
сверены по полным текстам; подробности --- в \texttt{sources/results.md}.

\section{Что получено}
Всюду предполагается оценка~\eqref{eq:bound}. Результаты собраны в таблице.

\begin{center}
\renewcommand{\arraystretch}{1.35}
\begin{tabular}{|>{\raggedright\arraybackslash}p{0.43\textwidth}|>{\raggedright\arraybackslash}p{0.49\textwidth}|}
\hline
\textbf{Условие на $a$} & \textbf{Вывод}\\\hline
$a\in L_\infty$ & обратный оператор компактен, верна альтернатива Фредгольма; базисность
не следует\\\hline
$a(-x)=a(x)$ & собственные функции образуют базис Рисса, присоединённых функций
нет\\\hline
$a\in W^2_\infty$, $0<m\le|a|\le M$ & корневые функции полны и образуют базис Рисса со
скобками\\\hline
то же и $\bigl|\frac{a'(x)}{a(x)}+\frac{a'(-x)}{a(-x)}\bigr|\le\eta$,
$160(1+\sqrt\kappa)\eta<\pi(1-\sqrt\kappa)$ & тот же вывод с элементарным
доказательством\\\hline
$a=4$ на $(0,1)$, $a=0$ на $(-1,0)$ & собственных значений нет\\\hline
$a(x)=\max(x,0)$ & нет базиса Рисса из собственных функций\\\hline
\end{tabular}
\end{center}

\textbf{Весь класс.} Так как $(aR)^2$ есть умножение на $a(x)a(-x)$, оператор $I+aR$
обратим, и обратный к нему выписывается явно. Поэтому $L_a^{-1}$ --- произведение
обратного оператора Дирихле и ограниченного оператора, то есть компактный оператор.
Спектр дискретен, но о базисности это ничего не говорит. При $\kappa=1$ утверждение
неверно: для $a\equiv1$ ядро $L_a$ бесконечномерно.

\textbf{Чётный коэффициент.} Чётные и нечётные функции образуют инвариантные
подпространства, и на них $L_a$ действует как $-(1\pm a)\,d^2/dx^2$. В весовом скалярном
произведении этот оператор самосопряжён, а вес ограничен сверху и снизу. Отсюда базис
Рисса и отсутствие присоединённых функций. Гладкость $a$ не нужна. Оператор $I+aR$
симметричен только для чётных $a$, поэтому на остальные коэффициенты этот приём не
переносится.

\textbf{Коэффициент постоянного знака.} Умножение на вес $\sqrt{|a|}$ превращает $L_a$ в
оператор с чётным коэффициентом $\pm\sqrt{a(x)a(-x)}$ плюс слагаемое первого порядка.
Первая часть диагонализуется предыдущим результатом, её собственные значения растут
квадратично, а слагаемое первого порядка подчинено ей с показателем $1/2$. Это в точности
условия теоремы Маркуса--Мацаева, и она даёт базис со скобками. Пример:
$a(x)=\tfrac12+\tfrac x4$.

\textbf{Коэффициент, близкий к чётному.} Если мала производная нечётной части функции
$\ln|a|$, то вместо теоремы Маркуса--Мацаева работает метод подобных операторов
(ср.~\cite{S5}): спектр разрезается на далеко отстоящие группы, и возмущённый оператор
приводится к блочно-диагональному виду сжимающим отображением. Контурные интегралы не
нужны. Пример: $a(x)=\tfrac12+\tfrac x{1000}$.

\textbf{Два примера.} В обоих $a(x)a(-x)=0$, то есть оценка~\eqref{eq:bound} выполнена
наилучшим образом, а половины интервала связаны только в одну сторону. Для разрывного
коэффициента задача сводится к задаче Коши и собственных значений нет вовсе. Для
липшицева $a(x)=\max(x,0)$ собственные значения лежат далеко от положительной полуоси,
вдоль которой есть почти собственные функции; такое расположение несовместимо с базисом
Рисса из собственных функций. Значит, одной оценки~\eqref{eq:bound} и обратимости
матрицы редукции для базисности недостаточно.

\textbf{Контрольные случаи.} При $a=0$ получается задача Дирихле. При $a=c$ собственные
функции --- $\cos((n+\tfrac12)\pi x)$ и $\sin(m\pi x)$ с собственными значениями
$(1+c)(n+\tfrac12)^2\pi^2$ и $(1-c)m^2\pi^2$. Коэффициент $a(x)=x/2$ удовлетворяет
оценке~\eqref{eq:bound} с $\kappa=\tfrac14$, но меняет знак, и ни один из результатов
таблицы к нему не применим.

\section{Что проверено и что открыто}
Все утверждения раздела~3 формализованы в Lean~4 с библиотекой Mathlib: библиотека
\texttt{Involution/} не содержит \texttt{sorry} и новых аксиом. Формализована и теорема
Маркуса--Мацаева в нужном частном случае: самосопряжённый положительный оператор,
показатель $1/2$. Теоремы сформулированы для области из функций с двумя слабыми
производными в $L_2$ и нулевыми значениями на концах; её совпадение с
$H^2\cap H_0^1$, где $H_0^1$ определено замыканием гладких функций, тоже формализовано.

Часть доказательств получена от сервиса Aristotle по формулировкам проекта: постоянство
функции с нулевой слабой производной, полнота системы косинусов и синусов, оценки
мультипликаторов Шура, леммы о разрезах спектра, лемма о весе для коэффициента, близкого
к чётному, теорема Маркуса--Мацаева по статье~\cite{S4} и описание $H_0^1$ через
нулевые значения на концах. Формулировки в ответах не
изменены, всё собрано локально. Остальное доказано локально; распределение по модулям
записано в \texttt{blueprint.md}.

Открыты следующие вопросы.
\begin{enumerate}
\item Образуют ли корневые функции обычный базис Рисса для коэффициента постоянного
знака? Для этого нужны равномерные константы Рисса внутри групп; они не оценены.
\item Каков размер групп в базисе со скобками?
\item Что происходит для коэффициентов с нулями и сменой знака, например $a(x)=x/2$?
Для нечётного $a$ формализовано тождество $UL_a=(1+ia)(-D^2)U$ с унитарным
$U=P_++iP_-$, где $P_\pm$ --- проекторы на чётные и нечётные функции; базисность для
задачи $-Z''=\lambda Z/(1+ia)$ не исследована.
\item Есть ли для $a(x)=\max(x,0)$ базис со скобками или базис с присоединёнными
функциями?
\end{enumerate}

\begin{thebibliography}{99}
\bibitem{S1}
Natalia P. Bondarenko. \emph{Inverse spectral problems for functional-differential operators with involution}.
arXiv:2107.12209v1. \url{https://arxiv.org/abs/2107.12209v1}.
\bibitem{S2}
Asylzat Kopzhassarova; Abdizhakhan Sarsenbi. \emph{Basis Properties of Eigenfunctions of Second-Order Differential Operators with Involution}.
2012, Article 576843. \url{https://doi.org/10.1155/2012/576843}.
\bibitem{S3}
А. А. Шкаликов. \emph{Возмущения самосопряженных и нормальных операторов с дискретным спектром}.
Успехи матем. наук 71:5(431) (2016), 113--174. \url{https://doi.org/10.4213/rm9740}.
\bibitem{S4}
Christian Wyss. \emph{Riesz Bases for p-Subordinate Perturbations of Normal Operators}.
arXiv:0906.2631v1. \url{https://arxiv.org/abs/0906.2631v1}.
\bibitem{S5}
Dmitry M. Polyakov. \emph{On the Bari basis property for even-order differential operators with involution}.
Tamkang J. Math. 54(4), 339--351 (2023); online 2022. \url{https://doi.org/10.5556/j.tkjm.54.2023.4899}.
\bibitem{S6}
Abdissalam Sarsenbi; Abdizhahan Sarsenbi. \emph{On Eigenfunctions of the Boundary Value Problems for Second Order Differential Equations with Involution}.
Symmetry 13(10), 1972 (2021). \url{https://doi.org/10.3390/sym13101972}.
\bibitem{S7}
Anton A. Lunyov; Mark M. Malamud. \emph{On the completeness and Riesz basis property of root subspaces of boundary value problems for first order systems and applications}.
arXiv:1401.2574v1. \url{https://arxiv.org/abs/1401.2574v1}.
\bibitem{S8}
Anton A. Lunyov; Mark M. Malamud. \emph{On the Riesz basis property of root vectors system for 2 × 2 Dirac type operators}.
arXiv:1504.04954v1. \url{https://arxiv.org/abs/1504.04954v1}.
\bibitem{S9}
Anton A. Lunyov; Mark M. Malamud. \emph{Stability of spectral characteristics and Bari basis property of boundary value problems for 2 × 2 Dirac type systems}.
arXiv:2012.11170v1. \url{https://arxiv.org/abs/2012.11170v1}.
\bibitem{S10}
Mark M. Malamud; Leonid L. Oridoroga. \emph{On the completeness of the system of root vectors for first-order systems}.
arXiv:1109.3683v1. \url{https://arxiv.org/abs/1109.3683v1}.
\bibitem{S11}
Boris Mityagin; Petr Siegl. \emph{Local form-subordination condition and Riesz basisness of root systems}.
arXiv:1608.00224v1. \url{https://arxiv.org/abs/1608.00224v1}.
\bibitem{S12}
Boris Mityagin; Petr Siegl. \emph{Local form subordination without a power decay and a criterion of Riesz basesness}.
arXiv:2411.10553v1. \url{https://arxiv.org/abs/2411.10553v1}.
\bibitem{S13}
Plamen Djakov; Boris Mityagin. \emph{Criteria for existence of Riesz bases consisting of root functions of Hill and 1D Dirac operators}.
arXiv:1106.5774v1. \url{https://arxiv.org/abs/1106.5774v1}.
\bibitem{S14}
Plamen Djakov; Boris Mityagin. \emph{Riesz bases consisting of root functions of 1D Dirac operators}.
arXiv:1108.4225v1. \url{https://arxiv.org/abs/1108.4225v1}.
\bibitem{S15}
Plamen Djakov; Boris Mityagin. \emph{Riesz basis property of Hill operators with potentials in weighted spaces}.
arXiv:1403.2973v1. \url{https://arxiv.org/abs/1403.2973v1}.
\bibitem{S16}
Plamen Djakov; Boris Mityagin. \emph{Equiconvergence of spectral decompositions of 1D Dirac operators with regular boundary conditions}.
arXiv:1108.0344v1. \url{https://arxiv.org/abs/1108.0344v1}.
\bibitem{S17}
Plamen Djakov; Boris Mityagin. \emph{Bari-Markus property for Riesz projections of 1D periodic Dirac operators}.
arXiv:0901.0856v1. \url{https://arxiv.org/abs/0901.0856v1}.
\bibitem{S18}
Plamen Djakov; Boris Mityagin. \emph{Bari-Markus property for Riesz projections of Hill operators with singular potentials}.
arXiv:0803.3170v1. \url{https://arxiv.org/abs/0803.3170v1}.
\bibitem{S19}
Hans Zwart. \emph{Riesz basis for strongly continuous groups}.
arXiv:0808.3447v1. \url{https://arxiv.org/abs/0808.3447v1}.
\bibitem{S20}
A. A. Shkalikov. \emph{Eigenvalue Asymptotics of Perturbed Self-adjoint Operators}.
arXiv:1202.5204v1. \url{https://arxiv.org/abs/1202.5204v1}.
\bibitem{S21}
A. A. Shkalikov; O. A. Veliev. \emph{On the Riesz Basis Property of the Eigen- and Associated Functions of Periodic and Antiperiodic Sturm-Liouville Problems}.
arXiv:0811.2337v1. \url{https://arxiv.org/abs/0811.2337v1}.
\end{thebibliography}
\end{document}
